\begin{abstract}
Classical Tucker models explain multiway data through multilinear low-rank
structure, but many measurement processes introduce nonlinear entry-level
effects such as saturation, response curves, and material-dependent
nonlinearities. We study a focused question: can one add nonlinear observation
capacity to Tucker while preserving its parameter-efficient low-rank
backbone?

We propose Nonlinear Tucker Decomposition with Polynomial Links (NTD-PL),
which applies a finite polynomial link to a shared Tucker latent tensor. The
model contains Tucker as a linear special case, acts as a compressed surrogate
for high-order Tucker interactions, and admits masked learning through
alternating gradient updates and a small ridge-regression subproblem for the
link coefficients. We give approximation and separation results, together
with a masked-entry generalization bound. Experiments on controlled tensors,
CAVE hyperspectral completion and reconstruction, and external HSI scenes
show that NTD-PL improves matched-rank Tucker when a shared nonlinear residual
is present, that the effect is not explained by modest rank increases or
post-hoc scalar calibration, and that boundary cases arise when Tucker already
leaves little coherent residual for the link to capture.
\end{abstract}
